\documentclass[11pt,letterpaper]{article}

% ==============================================================================
% PAQUETES DE CONFIGURACIÓN Y TIPOGRAFÍA
% ==============================================================================
\usepackage[utf8]{inputenc}
\usepackage[spanish,es-nodecimaldot,es-noshorthands]{babel}
\usepackage[margin=2.2cm,top=2.5cm,bottom=2.5cm]{geometry}
\usepackage{amsmath,amssymb,amsfonts,mathtools,bm}
\usepackage{microtype}
\usepackage{booktabs,array,tabularx}
\usepackage{fancyhdr}
\usepackage{lastpage}
\usepackage{xcolor}
\usepackage{tikz}
\usetikzlibrary{arrows.meta,patterns,calc,positioning,decorations.pathmorphing,shapes.geometric,babel}
\usepackage[skins,breakable]{tcolorbox}
\usepackage{hyperref}

% ==============================================================================
% PALETA DE COLORES INSTITUCIONAL Y PEDAGÓGICA
% ==============================================================================
\definecolor{uispurple}{RGB}{109, 40, 217}
\definecolor{uisblue}{RGB}{2, 132, 199}
\definecolor{uisgreen}{RGB}{5, 150, 105}
\definecolor{uisamber}{RGB}{217, 119, 6}
\definecolor{uisrose}{RGB}{225, 29, 72}
\definecolor{uisdark}{RGB}{15, 23, 42}
\definecolor{uisgray}{RGB}{100, 116, 139}
\definecolor{uislightbg}{RGB}{248, 250, 252}
\definecolor{uisboxborder}{RGB}{203, 213, 225}

% ==============================================================================
% HIPERVÍNCULOS
% ==============================================================================
\hypersetup{
    colorlinks=true,
    linkcolor=uispurple,
    citecolor=uisblue,
    urlcolor=uisblue,
    pdftitle={Equilibrio No Lineal y Metodo de Newton-Raphson - 2-DOF},
    pdfauthor={Prof. Orlando Arroyo - Universidad Industrial de Santander}
}

% ==============================================================================
% ENCABEZADOS Y PIES DE PÁGINA (FANCYHDR)
% ==============================================================================
\pagestyle{fancy}
\fancyhf{}
\renewcommand{\headrulewidth}{0.6pt}
\renewcommand{\footrulewidth}{0.4pt}
\fancyhead[L]{\small\textbf{\color{uispurple}Análisis No Lineal de Estructuras} | Posgrado en Ing. Estructural}
\fancyhead[R]{\small\color{uisgray}Universidad Industrial de Santander}
\fancyfoot[L]{\small\color{uisgray}Prof. Orlando Arroyo}
\fancyfoot[C]{\small\color{uisgray}Material Docente de Cátedra}
\fancyfoot[R]{\small\color{uisdark}\thepage\ / \pageref{LastPage}}

% ==============================================================================
% CAJAS DESTACADAS (TCOLORBOX)
% ==============================================================================
\newtcolorbox{conceptbox}[1][]{
    enhanced,
    breakable,
    colback=uislightbg,
    colframe=uispurple,
    coltitle=white,
    fonttitle=\bfseries\sffamily,
    title={#1},
    sharp corners=downhill,
    arc=4mm,
    boxrule=1.2pt,
    left=12pt,right=12pt,top=8pt,bottom=8pt,
    shadow={2pt}{-2pt}{0pt}{black!10}
}

\newtcolorbox{stepbox}[1][]{
    enhanced,
    breakable,
    colback=white,
    colframe=uisblue,
    coltitle=white,
    fonttitle=\bfseries\sffamily,
    title={#1},
    rounded corners,
    arc=3mm,
    boxrule=1.2pt,
    left=12pt,right=12pt,top=8pt,bottom=8pt
}

\newtcolorbox{alertbox}[1][]{
    enhanced,
    breakable,
    colback=uisrose!5,
    colframe=uisrose,
    coltitle=white,
    fonttitle=\bfseries\sffamily,
    title={#1},
    rounded corners,
    arc=3mm,
    boxrule=1.2pt,
    left=12pt,right=12pt,top=8pt,bottom=8pt
}

% ==============================================================================
% DOCUMENTO PRINCIPAL
% ==============================================================================
\begin{document}

% --- PORTADA / ENCABEZADO TITULAR ---
\begin{center}
    {\scshape\Large Universidad Industrial de Santander}\\[0.2cm]
    {\large Escuela de Ingeniería Civil --- Posgrado en Ingeniería Estructural}\\[0.35cm]
    {\color{uispurple}\hrule height 2pt}\vspace{0.4cm}
    {\huge\bfseries\color{uisdark} Formulación Teórica y Algorítmica del Equilibrio Estático No Lineal mediante el Método de Newton-Raphson}\\[0.25cm]
    {\Large\bfseries\color{uispurple} Sistema de Edificación Cortante de Dos Grados de Libertad (2-DOF)}\\[0.35cm]
    {\color{uispurple}\hrule height 1pt}\vspace{0.35cm}
    {\textbf{Autor:} Prof. Orlando Arroyo}\quad\textbullet\quad
    {\textbf{Asignatura:} Análisis No Lineal de Estructuras}\quad\textbullet\quad
    {\textbf{Versión:} Material Docente Consolidado (2026)}
\end{center}

\vspace{0.2cm}

\begin{abstract}
\noindent El presente documento constituye una guía formal, matemática y computacionalmente rigurosa, para comprender la resolución del equilibrio estático no lineal en estructuras discretizadas. Se desarrolla en detalle la deducción del vector de desequilibrio residual $\mathbf{r}(\mathbf{u})$, la linealización mediante series de Taylor, el ensamblaje explícito de la matriz de rigidez tangente o Jacobiana $\mathbf{K}_t$, y la solución incremental-iterativa según el algoritmo de Newton-Raphson. La teoría se aplica paso a paso sobre un modelo físico de edificio cortante de dos pisos (\textit{shear building} 2-DOF) con constitutivas bilineales elasto-plásticas, ilustrando de manera analítica y gráfica la evolución cinemática, la plastificación secuencial de entrepisos y la extinción cuadrática del desequilibrio de fuerzas.
\end{abstract}

\vspace{0.2cm}

% ==============================================================================
% SECCIÓN 1: FUNDAMENTOS DEL EQUILIBRIO NO LINEAL
% ==============================================================================
\section{Fundamentos del Equilibrio Estático No Lineal}

En mecánica computacional y análisis estructural inelástico, la respuesta resistente interna de la estructura deja de ser una función lineal de las coordenadas nodales. El equilibrio estático entre el vector de fuerzas externas aplicadas $\mathbf{P} \in \mathbb{R}^N$ y el vector de fuerzas resistentes internas $\mathbf{F}_{\text{int}}(\mathbf{u}) \in \mathbb{R}^N$ se expresa mediante la condición residual:
\begin{equation}
    \label{eq:residual}
    \mathbf{r}(\mathbf{u}) = \mathbf{P} - \mathbf{F}_{\text{int}}(\mathbf{u}) = \mathbf{0}
\end{equation}
donde $\mathbf{u} \in \mathbb{R}^N$ es el vector de desplazamientos generalizados y $\mathbf{r}(\mathbf{u})$ representa el \textbf{vector de desequilibrio o residuo dinámico}. A diferencia del análisis elástico lineal donde $\mathbf{F}_{\text{int}}(\mathbf{u}) = \mathbf{K}_0 \mathbf{u}$ permite una inversión directa $\mathbf{u} = \mathbf{K}_0^{-1} \mathbf{P}$, en el régimen no lineal la ecuación~(\ref{eq:residual}) representa un sistema algebraico acoplado de ecuaciones no lineales que debe resolverse mediante procesos iterativos.

\subsection{Linealización mediante Expansión en Series de Taylor}

Supóngase conocido un estado cinemático estimado en la iteración $i$, denotado como $\mathbf{u}^{(i)}$. Si el residuo $\mathbf{r}(\mathbf{u}^{(i)}) \neq \mathbf{0}$, se busca un incremento de desplazamiento correctivo $\delta \mathbf{u}^{(i)}$ tal que el nuevo estado $\mathbf{u}^{(i+1)} = \mathbf{u}^{(i)} + \delta \mathbf{u}^{(i)}$ satisfaga la condición de equilibrio:
\begin{equation}
    \mathbf{r}\left(\mathbf{u}^{(i)} + \delta \mathbf{u}^{(i)}\right) = \mathbf{0}
\end{equation}

Aproximando la función residual multivariable mediante una serie de Taylor de primer orden alrededor de $\mathbf{u}^{(i)}$:
\begin{equation}
    \mathbf{r}\left(\mathbf{u}^{(i)} + \delta \mathbf{u}^{(i)}\right) \approx \mathbf{r}\left(\mathbf{u}^{(i)}\right) + \left.\frac{\partial \mathbf{r}}{\partial \mathbf{u}}\right|_{\mathbf{u}^{(i)}} \delta \mathbf{u}^{(i)} + \mathcal{O}\left(\|\delta \mathbf{u}^{(i)}\|^2\right) = \mathbf{0}
\end{equation}

Puesto que las fuerzas externas aplicadas $\mathbf{P}$ en análisis controlado por carga no dependen cinemáticamente de $\mathbf{u}$ ($\partial \mathbf{P} / \partial \mathbf{u} = \mathbf{0}$), el Jacobiano del residuo resulta:
\begin{equation}
    \left.\frac{\partial \mathbf{r}}{\partial \mathbf{u}}\right|_{\mathbf{u}^{(i)}} = \frac{\partial \left(\mathbf{P} - \mathbf{F}_{\text{int}}\right)}{\partial \mathbf{u}} = -\left.\frac{\partial \mathbf{F}_{\text{int}}}{\partial \mathbf{u}}\right|_{\mathbf{u}^{(i)}} \equiv -\mathbf{K}_t^{(i)}
\end{equation}
donde $\mathbf{K}_t^{(i)} \in \mathbb{R}^{N \times N}$ es la \textbf{matriz de rigidez tangente} del sistema en el estado $\mathbf{u}^{(i)}$. Sustituyendo esta definición, se obtiene el sistema fundamental de ecuaciones de corrección de Newton-Raphson:

\begin{conceptbox}[Ecuación Fundamental de Corrección de Newton-Raphson]
\begin{equation}
    \label{eq:NR_system}
    \mathbf{K}_t^{(i)} \, \delta \mathbf{u}^{(i)} = \mathbf{r}^{(i)}
\end{equation}
con la correspondiente actualización cinemática:
\begin{equation}
    \label{eq:kinematic_update}
    \mathbf{u}^{(i+1)} = \mathbf{u}^{(i)} + \delta \mathbf{u}^{(i)}
\end{equation}
\end{conceptbox}

\subsection{Criterios de Convergencia}

La iteración continúa hasta que el estado estructural satisface tolerancias numéricas predefinidas. En la práctica de la ingeniería sismorresistente se evalúan dos criterios fundamentales:
\begin{enumerate}
    \item \textbf{Criterio de Fuerzas Desequilibradas (Residuo):}
    \begin{equation}
        \|\mathbf{r}^{(i)}\|_2 = \sqrt{\sum_{j=1}^N \left(r_j^{(i)}\right)^2} \le \epsilon_R
    \end{equation}
    donde $\epsilon_R$ es una tolerancia admisible de fuerza (e.g., $0.50\text{ kN}$ o $10^{-3} \|\mathbf{P}\|_2$).
    
    \item \textbf{Criterio de Incremento Cinemático:}
    \begin{equation}
        \|\delta \mathbf{u}^{(i)}\|_2 = \sqrt{\sum_{j=1}^N \left(\delta u_j^{(i)}\right)^2} \le \epsilon_U
    \end{equation}
\end{enumerate}

% ==============================================================================
% SECCIÓN 2: DESCRIPCIÓN DEL MODELO DE EDIFICIO 2-DOF
% ==============================================================================
\section{Cinemática y Modelo Constitutivo del Edificio Cortante (2-DOF)}

Considérese un pórtico o edificio cortante (\textit{shear building}) de dos niveles, caracterizado por losa de piso infinitamente rígida a flexión y deformación axial despreciable en las columnas. Los grados de libertad del sistema corresponden a los desplazamientos laterales horizontales de cada entrepiso:
\begin{equation}
    \mathbf{u} = \begin{Bmatrix} u_1 \\ u_2 \end{Bmatrix}
\end{equation}

\subsection{Cinemática de Deformación de Entrepiso (Derivas)}

El estado resistente de los elementos disipativos depende directamente de la deformación relativa o \textbf{deriva de entrepiso} $\Delta_j$, definida cinemáticamente como:
\begin{align}
    \Delta_1 &= u_1 - u_0 = u_1 \quad (\text{dado que la base está empotrada: } u_0 = 0) \label{eq:drift1} \\
    \Delta_2 &= u_2 - u_1 \label{eq:drift2}
\end{align}
En notación matricial, la cinemática se relaciona mediante la matriz de compatibilidad $\mathbf{B}$:
\begin{equation}
    \mathbf{\Delta} = \begin{Bmatrix} \Delta_1 \\ \Delta_2 \end{Bmatrix} = \begin{bmatrix} 1 & 0 \\ -1 & 1 \end{bmatrix} \begin{Bmatrix} u_1 \\ u_2 \end{Bmatrix} = \mathbf{B} \, \mathbf{u}
\end{equation}

\subsection{Ley Constitutiva Bilineal con Endurecimiento Cinemático}

Cada entrepiso $j \in \{1, 2\}$ actúa como un resorte cortante elasto-plástico bilineal gobernado por:
\begin{itemize}
    \item Rigidez elástica inicial: $k_{0j}$.
    \item Resistencia de fluencia al corte: $F_{yj}$.
    \item Desplazamiento de fluencia: $u_{yj} = F_{yj} / k_{0j}$.
    \item Razón de rigidez post-fluencia (\textit{hardening ratio}): $\alpha_j \in [0, 1]$.
    \item Rigidez tangente plástica: $k_{pj} = \alpha_j \, k_{0j}$.
\end{itemize}

El cortante de entrepiso resistente $V_j(\Delta_j)$ y la rigidez tangente instantánea $k_{tj}$ se evalúan analíticamente según:
\begin{equation}
    \label{eq:constitutive}
    V_j(\Delta_j) = 
    \begin{cases}
        k_{0j} \, \Delta_j, & \text{si } |\Delta_j| \le u_{yj} \quad (\text{\bfseries Régimen Elástico: } k_{tj} = k_{0j}) \\[0.2cm]
        \operatorname{sgn}(\Delta_j) \left[ F_{yj} + \alpha_j k_{0j} \left(|\Delta_j| - u_{yj}\right) \right], & \text{si } |\Delta_j| > u_{yj} \quad (\text{\bfseries Régimen Plástico: } k_{tj} = \alpha_j k_{0j})
    \end{cases}
\end{equation}

% --- FIGURA 1: MODELO DEL EDIFICIO Y DCL ---
\begin{figure}[htbp]
\centering
\begin{tikzpicture}[scale=0.92, >=Stealth]
    % --- SUBFIGURA A: EDIFICIO 2-DOF ---
    \begin{scope}[shift={(-4.0,0)}]
        % Terreno / Base
        \fill[pattern=north east lines, pattern color=uisgray!60] (-1.8,-0.35) rectangle (1.8,0);
        \draw[thick, line width=1.2pt, uisdark] (-1.8,0) -- (1.8,0);
        \node at (0,-0.60) {\small\textbf{Base Empotrada} ($u_0=0$)};
        
        % Columnas Piso 1
        \draw[line width=2.5pt, uisblue] (-1.0,0) -- (-1.0,2.35);
        \draw[line width=2.5pt, uisblue] (1.0,0) -- (1.0,2.35);
        % Resorte equivalente Piso 1
        \draw[decorate, decoration={zigzag, segment length=5.5, amplitude=4}, thick, uisblue] (0,0.3) -- (0,2.05);
        \node[right, font=\footnotesize, uisblue] at (0.22, 1.20) {$V_1(\Delta_1), \; k_{t1}$};
        
        % Diafragma Piso 1
        \node[draw=uisdark, line width=1.2pt, fill=uislightbg, rounded corners=2pt, inner sep=4pt, minimum width=2.6cm, font=\bfseries\footnotesize] (piso1) at (0,2.5) {Nivel 1 ($m_1$)};
        
        % Columnas Piso 2
        \draw[line width=2.5pt, uisrose] (-1.0,2.65) -- (-1.0,5.15);
        \draw[line width=2.5pt, uisrose] (1.0,2.65) -- (1.0,5.15);
        % Resorte equivalente Piso 2
        \draw[decorate, decoration={zigzag, segment length=5.5, amplitude=4}, thick, uisrose] (0,2.95) -- (0,4.85);
        \node[right, font=\footnotesize, uisrose] at (0.22, 3.90) {$V_2(\Delta_2), \; k_{t2}$};
        
        % Diafragma Piso 2
        \node[draw=uisdark, line width=1.2pt, fill=uislightbg, rounded corners=2pt, inner sep=4pt, minimum width=2.6cm, font=\bfseries\footnotesize] (piso2) at (0,5.3) {Nivel 2 ($m_2$)};
        
        % Cargas y Grados de Libertad Nivel 1
        \draw[->, line width=1.5pt, uisamber] (piso1.east) -- ++(1.1, 0) node[right, font=\small\bfseries] {$P_1$};
        \draw[->, line width=1.2pt, uispurple] (piso1.west) -- ++(-1.0, 0) node[left, font=\small\bfseries] {$u_1$};
        
        % Cargas y Grados de Libertad Nivel 2
        \draw[->, line width=1.5pt, uisamber] (piso2.east) -- ++(1.1, 0) node[right, font=\small\bfseries] {$P_2$};
        \draw[->, line width=1.2pt, uispurple] (piso2.west) -- ++(-1.0, 0) node[left, font=\small\bfseries] {$u_2$};

        \node[below, font=\bfseries\small, uisdark] at (0,-1.05) {(a) Modelo Cinemático 2-DOF};
    \end{scope}

    % --- SUBFIGURA B: DIAGRAMAS DE CUERPO LIBRE NODALES (DCL) ---
    \begin{scope}[shift={(4.0,0)}]
        % DCL Nivel 2
        \node[draw=uisdark, line width=1.2pt, fill=uislightbg, rounded corners=2pt, inner sep=4pt, minimum width=2.6cm, font=\bfseries\footnotesize] (nodo2) at (0,5.3) {Nodo 2 (Nivel 2)};
        
        \draw[->, line width=1.5pt, uisamber] (nodo2.east) -- ++(1.1, 0) node[right, font=\small\bfseries] {$P_2$};
        \draw[->, line width=1.5pt, uisrose] (nodo2.west) -- ++(-1.1, 0) node[left, font=\small\bfseries] {$V_2(\Delta_2)$};
        \node[above, font=\footnotesize, uisdark] at (0, 5.75) {$r_2 = P_2 - V_2$};

        % Separador entre pisos
        \draw[dashed, uisgray!50] (-2.5, 3.9) -- (2.5, 3.9);

        % DCL Nivel 1
        \node[draw=uisdark, line width=1.2pt, fill=uislightbg, rounded corners=2pt, inner sep=4pt, minimum width=2.6cm, font=\bfseries\footnotesize] (nodo1) at (0,2.5) {Nodo 1 (Nivel 1)};
        
        % Fuerzas en Nodo 1
        \draw[->, line width=1.5pt, uisamber] (nodo1.east) -- ++(1.1, 0) node[right, font=\small\bfseries] {$P_1$};
        \draw[->, line width=1.5pt, uisblue] (nodo1.west) -- ++(-1.1, 0) node[left, font=\small\bfseries] {$V_1(\Delta_1)$};
        \draw[->, line width=1.4pt, uisrose] (-0.8, 3.0) -- (0.8, 3.0) node[midway, above, font=\footnotesize\bfseries] {$V_2(\Delta_2)$ (corte sup.)};
        \node[below, font=\footnotesize, uisdark] at (0, 1.85) {$r_1 = P_1 - [V_1(\Delta_1) - V_2(\Delta_2)]$};

        \node[below, font=\bfseries\small, uisdark] at (0,-1.05) {(b) Diagramas de Cuerpo Libre (DCL)};
    \end{scope}
\end{tikzpicture}
\caption{Representación física del edificio cortante de dos pisos (2-DOF) y diagramas de cuerpo libre nodales.}
\label{fig:building_model}
\end{figure}

\subsection{Vector de Fuerzas Internas Nodales}

A partir de los diagramas de cuerpo libre nodales ilustrados en la Figura~\ref{fig:building_model}(b):
\begin{itemize}
    \item \textbf{Equilibrio en el Nodo 2 (Techo):} La única fuerza resistente que se opone a $P_2$ es el esfuerzo cortante generado por el resorte del segundo entrepiso:
    \begin{equation}
        F_{\text{int},2} = V_2(\Delta_2)
    \end{equation}
    
    \item \textbf{Equilibrio en el Nodo 1 (Piso intermedio):} El nodo recibe el cortante resistente del resorte inferior $V_1$ transmitido hacia la base, mientras que el resorte superior le transfiere en dirección contraria el cortante $V_2$:
    \begin{equation}
        F_{\text{int},1} = V_1(\Delta_1) - V_2(\Delta_2)
    \end{equation}
\end{itemize}

Por consiguiente, el vector global de fuerzas internas nodales está dado por:
\begin{equation}
    \label{eq:Fint_vector}
    \mathbf{F}_{\text{int}}(\mathbf{u}) = 
    \begin{Bmatrix} 
        F_{\text{int},1}(u_1, u_2) \\ 
        F_{\text{int},2}(u_1, u_2) 
    \end{Bmatrix} = 
    \begin{Bmatrix} 
        V_1(u_1) - V_2(u_2 - u_1) \\ 
        V_2(u_2 - u_1) 
    \end{Bmatrix}
\end{equation}

Y el vector de desequilibrio residual $\mathbf{r}(\mathbf{u})$ adopta la forma explícita:
\begin{equation}
    \label{eq:residual_vector}
    \mathbf{r}(\mathbf{u}) = \mathbf{P} - \mathbf{F}_{\text{int}}(\mathbf{u}) = 
    \begin{Bmatrix} 
        P_1 - \left[ V_1(u_1) - V_2(u_2 - u_1) \right] \\ 
        P_2 - V_2(u_2 - u_1) 
    \end{Bmatrix}
\end{equation}

% ==============================================================================
% SECCIÓN 3: ENSAMBLAJE DE LA MATRIZ DE RIGIDEZ TANGENTE
% ==============================================================================
\section{Ensamblaje Analítico de la Matriz de Rigidez Tangente \texorpdfstring{$\mathbf{K}_t$}{Kt}}

La matriz de rigidez tangente $\mathbf{K}_t$ se obtiene derivando formalmente el vector de fuerzas internas respecto a los grados de libertad nodales $\mathbf{u} = [u_1, u_2]^T$, aplicando la regla de la cadena multivariable:
\begin{equation}
    \mathbf{K}_t = \frac{\partial \mathbf{F}_{\text{int}}}{\partial \mathbf{u}} = 
    \begin{bmatrix}
        \dfrac{\partial F_{\text{int},1}}{\partial u_1} & \dfrac{\partial F_{\text{int},1}}{\partial u_2} \\[0.3cm]
        \dfrac{\partial F_{\text{int},2}}{\partial u_1} & \dfrac{\partial F_{\text{int},2}}{\partial u_2}
    \end{bmatrix}
\end{equation}

\subsection{Deducción Término a Término}

\begin{enumerate}
    \item \textbf{Término $K_{11} = \dfrac{\partial F_{\text{int},1}}{\partial u_1}$:}
    \begin{equation}
        \frac{\partial F_{\text{int},1}}{\partial u_1} = \frac{\partial \left(V_1 - V_2\right)}{\partial u_1} = \frac{\partial V_1}{\partial \Delta_1} \frac{\partial \Delta_1}{\partial u_1} - \frac{\partial V_2}{\partial \Delta_2} \frac{\partial \Delta_2}{\partial u_1}
    \end{equation}
    Puesto que $\frac{\partial \Delta_1}{\partial u_1} = 1$ y $\frac{\partial \Delta_2}{\partial u_1} = \frac{\partial (u_2 - u_1)}{\partial u_1} = -1$:
    \begin{equation}
        K_{11} = k_{t1} \cdot (1) - k_{t2} \cdot (-1) = k_{t1} + k_{t2}
    \end{equation}
    \textit{Significado físico:} Al mover unitariamente el Nivel 1 ($\delta u_1 = 1$), se deforma tanto el resorte inferior como el resorte superior, por lo cual ambas rigideces tangentes se suman directamente.

    \item \textbf{Término $K_{12} = \dfrac{\partial F_{\text{int},1}}{\partial u_2}$:}
    \begin{equation}
        K_{12} = \frac{\partial \left(V_1 - V_2\right)}{\partial u_2} = \underbrace{\frac{\partial V_1}{\partial u_2}}_{=0} - \frac{\partial V_2}{\partial \Delta_2} \frac{\partial \Delta_2}{\partial u_2} = -k_{t2} \cdot (1) = -k_{t2}
    \end{equation}

    \item \textbf{Término $K_{21} = \dfrac{\partial F_{\text{int},2}}{\partial u_1}$:}
    \begin{equation}
        K_{21} = \frac{\partial V_2}{\partial u_1} = \frac{\partial V_2}{\partial \Delta_2} \frac{\partial \Delta_2}{\partial u_1} = k_{t2} \cdot (-1) = -k_{t2}
    \end{equation}
    Nótese que $K_{12} = K_{21} = -k_{t2}$, garantizando la \textbf{simetría} de la matriz $\mathbf{K}_t$ en virtud del teorema de reciprocidad de Maxwell-Betti.

    \item \textbf{Término $K_{22} = \dfrac{\partial F_{\text{int},2}}{\partial u_2}$:}
    \begin{equation}
        K_{22} = \frac{\partial V_2}{\partial u_2} = \frac{\partial V_2}{\partial \Delta_2} \frac{\partial \Delta_2}{\partial u_2} = k_{t2} \cdot (1) = k_{t2}
    \end{equation}
    \textit{Significado físico:} Al desplazar unitariamente el Nivel 2 ($\delta u_2 = 1$), únicamente el resorte 2 experimenta deformación.
\end{enumerate}

\begin{conceptbox}[Estructura Matricial Tangente Ensamblada]
\begin{equation}
    \label{eq:Kt_assembled}
    \mathbf{K}_t = 
    \begin{bmatrix}
        k_{t1} + k_{t2} & -k_{t2} \\
        -k_{t2} & k_{t2}
    \end{bmatrix}
\end{equation}
\end{conceptbox}

\subsection{Determinante y Análisis de Estabilidad Mecánica}

El determinante de la matriz tangente resulta:
\begin{equation}
    \det(\mathbf{K}_t) = (k_{t1} + k_{t2})(k_{t2}) - (-k_{t2})(-k_{t2}) = k_{t1} k_{t2} + k_{t2}^2 - k_{t2}^2 = k_{t1} \cdot k_{t2}
\end{equation}

\begin{alertbox}[Pérdida de Capacidad Portante y Singularidad]
El determinante de $\mathbf{K}_t$ es exactamente igual al producto de las rigideces tangentes de entrepiso:
\begin{equation}
    \det(\mathbf{K}_t) = k_{t1} \cdot k_{t2}
\end{equation}
Si alguno de los entrepisos presenta comportamiento elasto-plástico perfecto ($\alpha_j = 0 \implies k_{tj} = 0$), el determinante se anula ($\det(\mathbf{K}_t) = 0$). La matriz se vuelve \textbf{singular}, el algoritmo de Newton-Raphson tradicional colapsa por división por cero y la estructura alcanza una cinemática de \textit{piso blando} o mecanismo plástico libre de carga.
\end{alertbox}

\subsection{Inversión Explícita y Flexibilidad Física de Corrección}

Al invertir analíticamente $\mathbf{K}_t$:
\begin{equation}
    \mathbf{K}_t^{-1} = \frac{1}{\det(\mathbf{K}_t)} 
    \begin{bmatrix} 
        K_{22} & -K_{12} \\ 
        -K_{21} & K_{11} 
    \end{bmatrix} = \frac{1}{k_{t1} k_{t2}} 
    \begin{bmatrix} 
        k_{t2} & k_{t2} \\ 
        k_{t2} & k_{t1} + k_{t2} 
    \end{bmatrix} = 
    \begin{bmatrix} 
        \dfrac{1}{k_{t1}} & \dfrac{1}{k_{t1}} \\[0.35cm] 
        \dfrac{1}{k_{t1}} & \dfrac{1}{k_{t1}} + \dfrac{1}{k_{t2}} 
    \end{bmatrix}
\end{equation}

Multiplicando $\mathbf{K}_t^{-1}$ por el vector residuo $\mathbf{r} = [r_1, r_2]^T$, se deduce la solución analítica de los incrementos:
\begin{align}
    \delta u_1 &= \frac{1}{k_{t1}} (r_1 + r_2) = \frac{r_{\text{basal}}}{k_{t1}} \label{eq:du1_physical} \\[0.25cm]
    \delta u_2 &= \frac{1}{k_{t1}} (r_1 + r_2) + \frac{1}{k_{t2}} r_2 = \delta u_1 + \frac{r_2}{k_{t2}} \label{eq:du2_physical}
\end{align}

\noindent\textbf{Interpretación Mecánica Fundamental:}
\begin{itemize}
    \item La corrección del Nivel 1 ($\delta u_1$) depende exclusivamente de la flexibilidad de la base ($1/k_{t1}$) y absorbe el \textbf{desequilibrio total de cortante sobre la base} ($r_1 + r_2$).
    \item La corrección del Nivel 2 ($\delta u_2$) arrastra el desplazamiento de la base ($\delta u_1$) y le suma la deformación interna producida por la flexibilidad del entrepiso superior ($1/k_{t2}$) sometido al desequilibrio de su propio nivel ($r_2$).
\end{itemize}

% --- FIGURA 2: CONSTITUTIVA Y NEWTON 1D ---
% --- FIGURA 2: CONSTITUTIVA Y NEWTON 1D ---
\begin{figure}[htbp]
\centering
\begin{tikzpicture}[scale=0.88, >=Stealth]
    % --- SUBFIGURA A: CURVA CONSTITUTIVA BILINEAL ---
    % --- SUBFIGURA A: CURVA CONSTITUTIVA BILINEAL ---
    \begin{scope}[shift={(-4.2,0)}]
        % Ejes
        \draw[->, thick, uisdark] (-0.3,0) -- (4.4,0) node[right, font=\footnotesize\bfseries] {$\Delta_j$ [mm]};
        \draw[->, thick, uisdark] (0,-0.3) -- (0,4.2) node[above, font=\footnotesize\bfseries] {$V_j$ [kN]};

        % Puntos clave
        \coordinate (O) at (0,0);
        \coordinate (Y) at (1.5,2.4);
        \coordinate (P) at (4.0,3.3);

        % Curva bilineal
        \draw[line width=2pt, uisgreen] (O) -- (Y);
        \draw[line width=2pt, uisrose] (Y) -- (P);

        % Líneas de proyección
        \draw[dashed, uisgray] (1.5,0) -- (Y) -- (0,2.4);
        \fill[uisdark] (Y) circle (2.5pt);
        \node[above left, font=\scriptsize, uisdark] at (Y) {Fluencia $(u_{yj}, F_{yj})$};

        % Etiquetas de ejes
        \node[below, font=\scriptsize] at (1.5,0) {$u_{yj}$};
        \node[left, font=\scriptsize] at (0,2.4) {$F_{yj}$};

        % Pendientes
        \draw[uisgreen, thick] (0.6,0.4) -- (1.1,0.4) -- (1.1,1.2);
        \node[right, font=\scriptsize\bfseries, uisgreen] at (1.1,0.75) {$k_{0j}$ (elástica)};

        \draw[uisrose, thick] (2.4,2.7) -- (3.4,2.7) -- (3.4,3.02);
        \node[above, font=\scriptsize\bfseries, uisrose] at (2.9,3.05) {$k_{pj} = \alpha_j k_{0j}$};

        \node[below, font=\bfseries\small, uisdark] at (2.0,-0.75) {(a) Relación Bilineal};
    \end{scope}

    % --- SUBFIGURA B: PROYECCIÓN NEWTON-RAPHSON ---
    \begin{scope}[shift={(4.2,0)}]
        % Ejes
        \draw[->, thick, uisdark] (-0.3,0) -- (4.6,0) node[right, font=\footnotesize\bfseries] {$u$ [mm]};
        \draw[->, thick, uisdark] (0,-0.3) -- (0,4.2) node[above, font=\footnotesize\bfseries] {$F$ [kN]};

        % Curva real
        \draw[line width=1.8pt, uisgray!50] (0,0) -- (1.2,1.8) -- (4.2,2.8);
        \node[above, font=\scriptsize, uisgray] at (3.8,2.75) {$F_{\text{int}}(u)$};

        % Carga objetivo P
        \draw[dashed, line width=1.2pt, uisamber] (0,3.3) -- (4.3,3.3) node[right, font=\scriptsize\bfseries] {Carga $P$};

        % Iteración 0: (0,0) -> Salto tangente elástico
        \fill[uisblue] (0,0) circle (2.2pt) node[below left, font=\scriptsize] {$u^{(0)}$};
        \draw[line width=1.2pt, uisblue, dashed] (0,0) -- (2.2, 3.3);
        \node[below, font=\scriptsize, uisblue] at (2.2, 0) {$u^{(1)}$};
        \draw[dotted, uisgray] (2.2,0) -- (2.2, 3.3);

        % Iteración 1: punto en curva
        \fill[uisrose] (2.2, 2.13) circle (2.2pt) node[right=3pt, font=\scriptsize] {$F_{\text{int}}(u^{(1)})$};
        
        % Residuo vertical
        \draw[line width=1.4pt, uisrose, <->] (2.2, 2.13) -- (2.2, 3.3) node[midway, left, font=\scriptsize\bfseries] {$r^{(1)}$};

        % Tangente plástica
        \draw[line width=1.2pt, uispurple] (2.2, 2.13) -- (3.9, 3.3);
        \fill[uispurple] (3.9, 3.3) circle (2.5pt) node[above left, font=\scriptsize\bfseries] {Solución $u^*$};
        \draw[dotted, uisgray] (3.9,0) -- (3.9, 3.3);
        \node[below, font=\scriptsize, uispurple] at (3.9, 0) {$u^{(2)}$};

        \node[below, font=\bfseries\small, uisdark] at (2.0,-0.75) {(b) Proyección Newton-Raphson};
    \end{scope}
\end{tikzpicture}
\caption{Relación constitutiva bilineal y esquema geométrico del algoritmo incremental de Newton-Raphson.}
\label{fig:constitutive_and_nr}
\end{figure}

% ==============================================================================
% SECCIÓN 4: CASO DE ESTUDIO NUMÉRICO RESUELTO PASO A PASO
% ==============================================================================
\section{Ejemplo de Aplicación Numérica Resuelto Paso a Paso}

Para consolidar la comprensión práctica del algoritmo, se resuelve a continuación un caso numérico idéntico al implementado en el simulador interactivo de clase.

\subsection{Parámetros de la Estructura y Solicitaciones}

\begin{itemize}
    \item \textbf{Propiedades del Entrepiso 1 (Base):}
    $k_{01} = 100.0\text{ kN/mm}$, $F_{y1} = 100.0\text{ kN}$, $\alpha_1 = 0.10$ ($10\%$).
    \[
        u_{y1} = \frac{100.0}{100.0} = 1.00\text{ mm}, \quad k_{p1} = 0.10 \times 100.0 = 10.0\text{ kN/mm}
    \]
    \item \textbf{Propiedades del Entrepiso 2 (Superior):}
    $k_{02} = 80.0\text{ kN/mm}$, $F_{y2} = 80.0\text{ kN}$, $\alpha_2 = 0.10$ ($10\%$).
    \[
        u_{y2} = \frac{80.0}{80.0} = 1.00\text{ mm}, \quad k_{p2} = 0.10 \times 80.0 = 8.0\text{ kN/mm}
    \]
    \item \textbf{Vector de Cargas Laterales Aplicadas:}
    \[
        \mathbf{P} = \begin{Bmatrix} P_1 \\ P_2 \end{Bmatrix} = \begin{Bmatrix} 60.0 \\ 110.0 \end{Bmatrix} \text{ kN}
    \]
    \textit{Observación previa:} El cortante total en la base es $P_1 + P_2 = 170.0\text{ kN} > F_{y1} = 100.0\text{ kN}$, y la fuerza en el segundo nivel $P_2 = 110.0\text{ kN} > F_{y2} = 80.0\text{ kN}$. Ambos pisos experimentarán plastificación forzosa.
    \item \textbf{Tolerancia de Convergencia:} $\epsilon_R = 0.50\text{ kN}$.
\end{itemize}

% ------------------------------------------------------------------------------
% ITERACIÓN 0
% ------------------------------------------------------------------------------
\subsection{Iteración \texorpdfstring{$i = 0$}{i = 0}}

\begin{stepbox}[Paso 0: Estado Inicial Elástico]
\begin{enumerate}
    \item \textbf{Vector cinemático inicial:} $\mathbf{u}^{(0)} = \begin{Bmatrix} 0.000 \\ 0.000 \end{Bmatrix}\text{ mm}$.
    \item \textbf{Derivas de entrepiso:} $\Delta_1 = u_1 = 0.000\text{ mm}$, $\Delta_2 = u_2 - u_1 = 0.000\text{ mm}$.
    \item \textbf{Cortantes resistentes:} $V_1 = 0.00\text{ kN}$, $V_2 = 0.00\text{ kN}$. Ambos pisos están en régimen \textbf{Elástico}:
    \[
        k_{t1}^{(0)} = k_{01} = 100.0\text{ kN/mm}, \quad k_{t2}^{(0)} = k_{02} = 80.0\text{ kN/mm}
    \]
    \item \textbf{Fuerzas internas nodales:}
    \[
        \mathbf{F}_{\text{int}}^{(0)} = \begin{Bmatrix} V_1 - V_2 \\ V_2 \end{Bmatrix} = \begin{Bmatrix} 0.00 \\ 0.00 \end{Bmatrix}\text{ kN}
    \]
    \item \textbf{Vector residuo de desequilibrio:}
    \[
        \mathbf{r}^{(0)} = \mathbf{P} - \mathbf{F}_{\text{int}}^{(0)} = \begin{Bmatrix} 60.00 - 0.00 \\ 110.00 - 0.00 \end{Bmatrix} = \begin{Bmatrix} +60.00 \\ +110.00 \end{Bmatrix}\text{ kN}
    \]
    \[
        \|\mathbf{r}^{(0)}\|_2 = \sqrt{(60.00)^2 + (110.00)^2} = 125.300\text{ kN} > 0.50\text{ kN} \quad (\text{\bfseries No converge})
    \]
    \item \textbf{Ensamblaje de la matriz tangente inicial $\mathbf{K}_t^{(0)}$:}
    \[
        \mathbf{K}_t^{(0)} = \begin{bmatrix} 100.0 + 80.0 & -80.0 \\ -80.0 & 80.0 \end{bmatrix} = \begin{bmatrix} 180.0 & -80.0 \\ -80.0 & 80.0 \end{bmatrix}\text{ kN/mm}
    \]
    \[
        \det(\mathbf{K}_t^{(0)}) = 100.0 \times 80.0 = 8000.0
    \]
    \item \textbf{Solución del sistema $\mathbf{K}_t^{(0)} \delta \mathbf{u}^{(0)} = \mathbf{r}^{(0)}$:}
    \[
        \begin{bmatrix} 180.0 & -80.0 \\ -80.0 & 80.0 \end{bmatrix} \begin{Bmatrix} \delta u_1^{(0)} \\ \delta u_2^{(0)} \end{Bmatrix} = \begin{Bmatrix} +60.00 \\ +110.00 \end{Bmatrix}
    \]
    Aplicando las ecuaciones analíticas~(\ref{eq:du1_physical})--(\ref{eq:du2_physical}):
    \[
        \delta u_1^{(0)} = \frac{r_1 + r_2}{k_{t1}} = \frac{60.00 + 110.00}{100.0} = \frac{170.00}{100.0} = \mathbf{+1.700\text{ mm}}
    \]
    \[
        \delta u_2^{(0)} = \delta u_1^{(0)} + \frac{r_2}{k_{t2}} = 1.700 + \frac{110.00}{80.0} = 1.700 + 1.375 = \mathbf{+3.075\text{ mm}}
    \]
    \item \textbf{Actualización cinemática para la siguiente iteración:}
    \[
        \mathbf{u}^{(1)} = \mathbf{u}^{(0)} + \delta \mathbf{u}^{(0)} = \begin{Bmatrix} 0.000 + 1.700 \\ 0.000 + 3.075 \end{Bmatrix} = \begin{Bmatrix} \mathbf{1.700} \\ \mathbf{3.075} \end{Bmatrix}\text{ mm}
    \]
\end{enumerate}
\end{stepbox}

% ------------------------------------------------------------------------------
% ITERACIÓN 1
% ------------------------------------------------------------------------------
\subsection{Iteración \texorpdfstring{$i = 1$}{i = 1}}

\begin{stepbox}[Paso 1: Transición Inelástica y Plastificación de Ambos Entrepisos]
\begin{enumerate}
    \item \textbf{Estado cinemático actual:} $\mathbf{u}^{(1)} = \begin{Bmatrix} 1.700 \\ 3.075 \end{Bmatrix}\text{ mm}$.
    \item \textbf{Evaluación de derivas de entrepiso:}
    \begin{align*}
        \Delta_1 &= u_1 = 1.700\text{ mm} > u_{y1} = 1.00\text{ mm} \quad \implies \quad \text{\bfseries\color{uisrose} ¡Piso 1 en Régimen Plástico!} \\
        \Delta_2 &= u_2 - u_1 = 1.375\text{ mm} > u_{y2} = 1.00\text{ mm} \quad \implies \quad \text{\bfseries\color{uisrose} ¡Piso 2 en Régimen Plástico!}
    \end{align*}
    \item \textbf{Cálculo de cortantes resistentes inelásticos:}
    \[
        V_1 = F_{y1} + k_{p1}(\Delta_1 - u_{y1}) = 100.0 + 10.0(1.700 - 1.00) = 100.0 + 7.00 = \mathbf{107.00\text{ kN}}
    \]
    \[
        V_2 = F_{y2} + k_{p2}(\Delta_2 - u_{y2}) = 80.0 + 8.0(1.375 - 1.00) = 80.0 + 3.00 = \mathbf{83.00\text{ kN}}
    \]
    Rigideces tangentes activas:
    \[
        k_{t1}^{(1)} = k_{p1} = \mathbf{10.0\text{ kN/mm}}, \quad k_{t2}^{(1)} = k_{p2} = \mathbf{8.0\text{ kN/mm}}
    \]
    \item \textbf{Fuerzas internas nodales:}
    \[
        F_{\text{int},1}^{(1)} = V_1 - V_2 = 107.00 - 83.00 = \mathbf{24.00\text{ kN}}
    \]
    \[
        F_{\text{int},2}^{(1)} = V_2 = \mathbf{83.00\text{ kN}}
    \]
    \item \textbf{Vector residuo de desequilibrio:}
    \[
        r_1^{(1)} = P_1 - F_{\text{int},1}^{(1)} = 60.00 - 24.00 = \mathbf{+36.00\text{ kN}}
    \]
    \[
        r_2^{(1)} = P_2 - F_{\text{int},2}^{(1)} = 110.00 - 83.00 = \mathbf{+27.00\text{ kN}}
    \]
    \[
        \|\mathbf{r}^{(1)}\|_2 = \sqrt{(36.00)^2 + (27.00)^2} = \sqrt{1296 + 729} = \sqrt{2025} = \mathbf{45.000\text{ kN}} > 0.50\text{ kN}
    \]
    \item \textbf{Ensamblaje de la nueva matriz tangente plástica $\mathbf{K}_t^{(1)}$:}
    \[
        \mathbf{K}_t^{(1)} = \begin{bmatrix} 10.0 + 8.0 & -8.0 \\ -8.0 & 8.0 \end{bmatrix} = \begin{bmatrix} \mathbf{18.0} & \mathbf{-8.0} \\ \mathbf{-8.0} & \mathbf{8.0} \end{bmatrix}\text{ kN/mm}
    \]
    \[
        \det(\mathbf{K}_t^{(1)}) = k_{t1}^{(1)} \cdot k_{t2}^{(1)} = 10.0 \times 8.0 = \mathbf{80.0}
    \]
    \textit{Comentario pedagógico:} Nótese la drástica reducción del determinante de $8000$ a $80$ (disminución del $99\%$), reflejo fiel de la pérdida de rigidez al plastificarse la estructura.
    \item \textbf{Solución del sistema de corrección $\mathbf{K}_t^{(1)} \delta \mathbf{u}^{(1)} = \mathbf{r}^{(1)}$:}
    \[
        \begin{bmatrix} 18.0 & -8.0 \\ -8.0 & 8.0 \end{bmatrix} \begin{Bmatrix} \delta u_1^{(1)} \\ \delta u_2^{(1)} \end{Bmatrix} = \begin{Bmatrix} +36.00 \\ +27.00 \end{Bmatrix}
    \]
    \[
        \delta u_1^{(1)} = \frac{r_1^{(1)} + r_2^{(1)}}{k_{t1}^{(1)}} = \frac{36.00 + 27.00}{10.0} = \frac{63.00}{10.0} = \mathbf{+6.300\text{ mm}}
    \]
    \[
        \delta u_2^{(1)} = \delta u_1^{(1)} + \frac{r_2^{(1)}}{k_{t2}^{(1)}} = 6.300 + \frac{27.00}{8.0} = 6.300 + 3.375 = \mathbf{+9.675\text{ mm}}
    \]
    \item \textbf{Actualización cinemática:}
    \[
        \mathbf{u}^{(2)} = \mathbf{u}^{(1)} + \delta \mathbf{u}^{(1)} = \begin{Bmatrix} 1.700 + 6.300 \\ 3.075 + 9.675 \end{Bmatrix} = \begin{Bmatrix} \mathbf{8.000} \\ \mathbf{12.750} \end{Bmatrix}\text{ mm}
    \]
\end{enumerate}
\end{stepbox}

% ------------------------------------------------------------------------------
% ITERACIÓN 2
% ------------------------------------------------------------------------------
\subsection{Iteración \texorpdfstring{$i = 2$}{i = 2}}

\begin{stepbox}[Paso 2: Evaluación en Estado Final y Verificación de Convergencia]
\begin{enumerate}
    \item \textbf{Estado cinemático actual:} $\mathbf{u}^{(2)} = \begin{Bmatrix} 8.000 \\ 12.750 \end{Bmatrix}\text{ mm}$.
    \item \textbf{Derivas de entrepiso:}
    \[
        \Delta_1 = u_1 = 8.000\text{ mm} > 1.00\text{ mm} \quad (\text{Plástico})
    \]
    \[
        \Delta_2 = u_2 - u_1 = 12.750 - 8.000 = 4.750\text{ mm} > 1.00\text{ mm} \quad (\text{Plástico})
    \]
    \item \textbf{Cortantes resistentes inelásticos evaluados:}
    \[
        V_1 = 100.0 + 10.0(8.000 - 1.00) = 100.0 + 70.00 = \mathbf{170.00\text{ kN}}
    \]
    \[
        V_2 = 80.0 + 8.0(4.750 - 1.00) = 80.0 + 30.00 = \mathbf{110.00\text{ kN}}
    \]
    \item \textbf{Fuerzas internas nodales:}
    \[
        F_{\text{int},1}^{(2)} = V_1 - V_2 = 170.00 - 110.00 = \mathbf{60.00\text{ kN}}
    \]
    \[
        F_{\text{int},2}^{(2)} = V_2 = \mathbf{110.00\text{ kN}}
    \]
    \item \textbf{Vector residuo de desequilibrio:}
    \[
        r_1^{(2)} = P_1 - F_{\text{int},1}^{(2)} = 60.00 - 60.00 = \mathbf{0.00\text{ kN}}
    \]
    \[
        r_2^{(2)} = P_2 - F_{\text{int},2}^{(2)} = 110.00 - 110.00 = \mathbf{0.00\text{ kN}}
    \]
    \[
        \|\mathbf{r}^{(2)}\|_2 = \sqrt{0^2 + 0^2} = \mathbf{0.000\text{ kN}} \le \epsilon_R = 0.50\text{ kN}
    \]
    \item \textbf{\color{uisgreen} ¡Convergencia Estricta Alcanzada!} El sistema se encuentra en equilibrio estático exacto en el estado $\mathbf{u}^* = [8.000, 12.750]^T\text{ mm}$.
\end{enumerate}
\end{stepbox}

\begin{alertbox}[Propiedad de Convergencia en Modelos Lineales por Tramos]
Nótese que en este sistema bilineal, una vez que el algoritmo identificó correctamente el régimen activo de ambos pisos (ambos en rama plástica con rigideces $k_{p1}$ y $k_{p2}$), el método de Newton-Raphson alcanzó la solución exacta en \textbf{exactamente una sola iteración correctiva adicional}. Esto ocurre porque sobre cada rama el modelo es afín (lineal por tramos), por lo que la aproximación de Taylor de primer orden resulta exacta.
\end{alertbox}

\subsection{Resumen Comparativo de Iteraciones}

En el Cuadro~\ref{tab:results_summary} se sintetiza la evolución cuantitativa del algoritmo Newton-Raphson estándar frente al comportamiento del método Newton-Raphson modificado (donde la matriz de rigidez se congela en la elástica inicial $\mathbf{K}_0$).

\begin{table}[htbp]
\centering
\footnotesize
\setlength{\tabcolsep}{2.5pt}
\renewcommand{\arraystretch}{1.25}
\caption{Historial de convergencia numérico detallado del sistema 2-DOF.}
\label{tab:results_summary}
\begin{tabular*}{\textwidth}{@{\extracolsep{\fill}} c c c c c c c c c l @{}}
\toprule
\textbf{Iter $i$} & $u_1\text{ [mm]}$ & $u_2\text{ [mm]}$ & $\Delta_1\text{ [mm]}$ & $\Delta_2\text{ [mm]}$ & $V_1\text{ [kN]}$ & $V_2\text{ [kN]}$ & $\|\mathbf{r}\|_2\text{ [kN]}$ & $\|\delta \mathbf{u}\|_2\text{ [mm]}$ & \textbf{Estado del Sistema} \\
\midrule
\multicolumn{10}{l}{\textbf{Método Newton-Raphson Estándar (Actualización Completa de $\mathbf{K}_t$):}} \\
0 & 0.000 & 0.000 & 0.000 & 0.000 & 0.0 & 0.0 & 125.30 & 3.514 & Elástico / Elástico \\
1 & 1.700 & 3.075 & 1.700 & 1.375 & 107.0 & 83.0 & 45.00 & 11.545 & Plástico / Plástico \\
2 & 8.000 & 12.750 & 8.000 & 4.750 & 170.0 & 110.0 & \textbf{0.00} & \textbf{0.000} & \textbf{Convergido (2 iter)} \\
\midrule
\multicolumn{10}{l}{\textbf{Método Newton-Raphson Modificado (Rigidez Constante $\mathbf{K}_0 = \mathbf{K}_t^{(0)}$):}} \\
0 & 0.000 & 0.000 & 0.000 & 0.000 & 0.0 & 0.0 & 125.30 & 3.514 & Elástico / Elástico \\
1 & 1.700 & 3.075 & 1.700 & 1.375 & 107.0 & 83.0 & 45.00 & 1.263 & Iterando... \\
2 & 2.330 & 4.218 & 2.330 & 1.888 & 113.3 & 87.1 & 37.80 & 1.061 & Iterando... \\
3 & 2.859 & 5.177 & 2.859 & 2.318 & 118.6 & 90.5 & 31.75 & 0.891 & Iterando... \\
$\vdots$ & $\vdots$ & $\vdots$ & $\vdots$ & $\vdots$ & $\vdots$ & $\vdots$ & $\vdots$ & $\vdots$ & $\vdots$ \\
28 & 7.965 & 12.687 & 7.965 & 4.722 & 169.7 & 109.8 & 0.49 & 0.014 & \textbf{Convergido (28 iter)} \\
\bottomrule
\end{tabular*}
\end{table}

% ==============================================================================
% SECCIÓN 5: COMPARACIÓN DE MÉTODOS Y CONSIDERACIONES INGENIERILES
% ==============================================================================
\section{Aspectos Teóricos y Prácticos en Ingeniería Estructural}

\subsection{Newton-Raphson Estándar vs. Modificado}

\begin{itemize}
    \item \textbf{Newton-Raphson Estándar (Consistente):}
    \begin{itemize}
        \item \textit{Ventaja:} Tasa de convergencia cuadrática en la proximidad de la solución ($\|\mathbf{r}^{(i+1)}\| \propto \|\mathbf{r}^{(i)}\|^2$). En el ejemplo, requirió únicamente 2 iteraciones.
        \item \textit{Desventaja:} Requiere formar, ensamblar y factorizar la matriz $\mathbf{K}_t$ en cada iteración, lo que en modelos de elementos finitos con cientos de miles de grados de libertad resulta computacionalmente muy demandante.
    \end{itemize}

    \item \textbf{Newton-Raphson Modificado (Initial Stiffness):}
    \begin{itemize}
        \item \textit{Ventaja:} La matriz $\mathbf{K}_0$ se factoriza una sola vez al inicio del análisis. Las iteraciones subsecuentes consisten en sustituciones progresivas y regresivas de bajo costo computacional.
        \item \textit{Desventaja:} Tasa de convergencia puramente lineal. Cuando la estructura experimenta plastificación pronunciada ($\alpha \ll 1$), la rigidez inicial elástica sobreestima enormemente la rigidez real, provocando que el algoritmo requiera decenas de iteraciones (28 pasos en nuestro ejemplo) o incluso diverja.
    \end{itemize}
\end{itemize}

\subsection{Sensibilidad a la Pérdida de Capacidad Portante}

Cuando un sistema estructural experimenta degradación de resistencia (\textit{strain softening}, donde $\alpha < 0$), los términos diagonales de $\mathbf{K}_t$ se reducen hasta que $\det(\mathbf{K}_t) \le 0$. En ese instante:
\begin{enumerate}
    \item El método controlado por carga de Newton-Raphson es incapaz de superar el punto límite de carga (\textit{limit point}).
    \item Se hace mandatorio migrar a algoritmos avanzados de \textbf{Control por Desplazamiento} o técnicas de \textbf{Longitud de Arco} (\textit{Arc-Length Methods} de Riks-Wempner y Crisfield), donde el factor de carga $\lambda$ pasa a ser una variable cinemática desconocida del sistema.
\end{enumerate}

% ==============================================================================
% CONCLUSIONES
% ==============================================================================
\section{Conclusiones Pedagógicas}

\begin{enumerate}
    \item El equilibrio estático en estructuras no lineales no se satisface de forma directa, sino mediante la anulación progresiva de un vector residual de fuerzas $\mathbf{r}(\mathbf{u}) = \mathbf{P} - \mathbf{F}_{\text{int}}(\mathbf{u})$.
    \item La matriz de rigidez tangente $\mathbf{K}_t$ no es una constante del material, sino el operador Jacobiano instantáneo de las fuerzas internas nodales respecto a los desplazamientos.
    \item En sistemas de edificación cortante, el determinante de $\mathbf{K}_t$ se reduce al producto directo de las rigideces tangentes de entrepiso ($\prod k_{tj}$), lo que permite visualizar físicamente cómo la plastificación de un solo entrepiso disminuye en órdenes de magnitud la estabilidad global del edificio.
    \item La solución explícita de las correcciones cinemáticas desacopla la física del problema: el primer piso absorbe el desequilibrio de corte basal acumulado ($r_1 + r_2$), mientras que los pisos superiores acumulan dicho desplazamiento y añaden su propia flexibilidad relativa.
\end{enumerate}

% ==============================================================================
% REFERENCIAS
% ==============================================================================
\begin{thebibliography}{9}
    \bibitem{chopra} Chopra, A. K. (2020). \textit{Dynamics of Structures: Theory and Applications to Earthquake Engineering}. 5th Edition, Pearson.
    \bibitem{crisfield} Crisfield, M. A. (1991). \textit{Non-linear Finite Element Analysis of Solids and Structures: Essentials}. John Wiley \& Sons.
    \bibitem{bathe} Bathe, K. J. (2006). \textit{Finite Element Procedures}. Prentice Hall, Pearson Education.
    \bibitem{arroyo} Arroyo, O. (2026). \textit{Notas de Clase y Laboratorios Computacionales de Análisis No Lineal de Estructuras}. Universidad Industrial de Santander. Repositorio en línea: \url{https://odarroyo.github.io/analisis_no_lineal/}.
\end{thebibliography}

\end{document}
